
\section{Introduction}
\label{sec:intro}

Many tasks in scientific and engineering applications can be framed as \emph{bandit
optimisation} problems, where we need to sequentially evaluate a noisy
black-box function $\func:\Xcal\rightarrow\RR$ with the goal of finding its optimum.
Some applications include hyper-parameter tuning in machine learning~\citep{%
snoek12practicalBO,hutter11smac}, optimal policy search~\citep{lizotte07automaticGait,
martinez07robotplanning} and scientific experiments~\citep{parkinson06wmap3,gonzalez14gene}.
Typically, in such applications, each function evaluation is expensive and historically
the bandit literature has focused on developing methods for finding the optimum
while keeping the number of evaluations to $\func$ at a minimum.

However, with increasingly expensive experiments, conventional methods have become
infeasible as a significant cost needs to be expended before we can learn anything
about $\func$.
As a result, \emph{multi-fidelity} optimisation methods have recently gained attention.
As the name suggests, these methods assume that we have access to lower fidelity
approximate functions to $\func$ which can be evaluated instead of $\func$.
The lower the fidelity, the cheaper the evaluation, but it provides less accurate
information about $\func$.
For example, when optimising the configuration of an expensive real world robot,
its performance can be approximated using cheap computer simulations.
The goal is to use the cheap approximations to guide search for the optimum of $\func$,
and reduce the overall cost of optimisation.
However, most multi-fidelity work assume only a finite number of approximations to
$\func$.
In this paper, we study multi-fidelity optimisation when there is access to a continuous
spectrum of approximations.



To motivate this,
consider tuning a classification algorithm over a space of hyper-parameters $\Xcal$
by maximising a validation set accuracy. 
The algorithm is to be trained using $\Nmax$
data points via an iterative algorithm for $\Tmax$ iterations.
However, we wish to use fewer training points $N<\Nmax$ and/or fewer iterations
$T<\Tmax$ to approximate the validation accuracy.
We can view validation accuracy as a function
$g:[1,\Nmax]\times[1,\Tmax]\times\Xcal\rightarrow\RR$ where evaluating $g(N,T,x)$
requires training the algorithm with $N$ points for $T$ iterations with the
hyper-parameters $x$.
If the training complexity of the algorithm is quadratic in data size and linear in the
number of iterations, then the cost of this evaluation is $\lambda(N,T) = \bigO(N^2T)$.
Our goal is to find the optimum when $N=\Nmax$, and $T=\Tmax$,
% i.e.  $\xopt = \argmax_{x\in\Xcal} g(\Nmax,\Tmax, x)$.
i.e.  we wish to maximise $\func(x) = g(\Nmax,\Tmax, x)$.

In this setting, the approximations are best viewed as coming from a continuous $2$
dimensional \emph{fidelity space}, namely $[1, \Nmax]\times[1,\Tmax]$.
% Since querying $g$ at $N,T$ values less than $\Nmax,\Tmax$ costs less, we are
% tempted to use such queries in the hope that we may
Since querying $g(N,T,\cdot)$ at $N,T$ less than $\Nmax,\Tmax$ costs less, we are
tempted to use such queries in the hope that we may
learn about $g(\Nmax, \Tmax,\cdot)$ cheaply and consequently
optimise it using less overall cost.
This is, in fact, the case with many machine learning algorithms where cross
validation curves tend to vary smoothly with data set size and number of iterations.
Then we may use cheap low fidelity experiments with small $(N, T)$ to discard bad regions
of the hyper-parameter space and deploy expensive high fidelity experiments with large
$(N,T)$ only in a small but promising region of $\Xcal$.
Our proposed algorithm exhibits precisely this behaviour. 

Continuous approximations also arise in
simulation studies -- where simulations can be carried out at varying levels of granularity,
on-line advertising -- where the display of an ad can be controlled by various continuous
parameters such as display time or target audience, and various other
experiment design tasks.
In fact, in many papers on multi-fidelity optimisation, the finite approximations were
obtained by discretising
a continuous space~\citep{kandasamy16mfbo,huang06mfKriging}.
Here, we study a \bayo{} technique that is directly designed for continuous fidelity
spaces and could potentially be applied to more general spaces.
\textbf{Our contributions are as follows.}
\vspace{-0.09in}
\begin{enumerate}
\item A novel setting and model for multi-fidelity optimisation with continuous
approximations using Gaussian process (GP) assumptions.
We develop a novel algorithm, \boca, for this setting.
\vspace{-0.09in}
\item A theoretical analysis characterising the behaviour and regret bound for \boca.
\vspace{-0.09in}
\item Empirically we demonstrate that \bocas outperforms other single and multi-fidelity
  methods on a series of synthetic problems, and real problems in hyper-parameter tuning
  and computational astrophysics.
%   Our python implementation will be made publicly 
\end{enumerate}


\vspace{-0.07in}
\subsection*{Related Work}
\vspace{-0.07in}

Bayesian optimisation (\bayo), refers to a suite of techniques for bandit optimisation
which use a prior belief distribution for $\func$.
While there are several techniques for
BO~\citep{mockus94bo,jones98expensive,thompson33sampling,lobato14pes},
our work will build on the 
Gaussian process upper confidence bound (\gpucb) algorithm
of~\citet{srinivas10gpbandits}.
\gpucbs models $\func$ as a GP and uses upper confidence bound (UCB)~\citep{auer03ucb}
techniques to determine the next point for evaluation.

\bayo{} techniques have been used in developing multi-fidelity optimisation methods
in various applications such as hyper-parameter tuning and industrial
design~\citep{huang06mfKriging,swersky2013multi,klein2015towards,forrester07cokriging,%
poloczek2016multi}.
However, these methods are either problem specific and/or only use a finite
number of fidelities.
Further, none of them come with theoretical underpinnings.
Recent work has studied multi-fidelity methods for specific problems 
such as hyper-parameter tuning, active learning and
reinforcement learning%
~\citep{agarwal2011oracle,sabharwal2015selecting,cutler14mfsim,zhang15weakAndStrong,%
li2016hyperband}.
While some of the above tasks can be framed as optimisation problems, the methods
themselves are specific to the problem considered.
Our method is more general as it applies to any bandit optimisation task.


Perhaps the closest work to us is that of~\citet{kandasamy16mfbo,kandasamy16mfbandit}
who developed \mfgpucbones assuming a finite number of approximations to $\func$.
% This line of work was important as it provided the first theoretical guarantees for
% multi-fidelity optimisation.
While this line of work was the first to provide theoretical guarantees for
multi-fidelity optimisation, it has two important shortcomings.
First, they make strong assumptions, particularly a uniform bound on the
difference between the expensive function and an approximation.
This does not allow for instances where an approximation might be good at certain regions
but not at the other.
In contrast our treatment between the fidelities is probabilistic which is robust
to such cases.
Second, their model does not allow sharing information between approximations; each
approximation is treated independently.
In addition to not uti
Not only is this wasteful because lower fidelities can provide useful information about
higher fidelities, it also means that the algorithm might perform poorly if the fidelities
are not designed properly.
We demonstrate this with an experiment in Section~\ref{sec:experiments}.
On the other hand, our model allows sharing information across the fidelity space in a
natural way.
In addition, we can also handle continuous approximations whereas their method is strictly
for a finite number of approximations.
That said, \mfgpucbtwos inherits a key intuition from \mfgpucbone,
% \todo{we don't need ``-I'' anymore right? Not reading the intro part too closely yet}
% first select the next point for evaluation using an upper confidence bound strategy;
which is to choose a fidelity only if we have controlled the uncertainty at all lower
fidelities.
% However, in addition to dealing with continuous approximations, we also differ from%
% ~\citet{kandasamy16mfbo} fundamentally in our model assumptions.
Besides this, the mechanics
of our algorithm and the analysis techniques are considerably different.
As we proceed, we will draw further comparisons to~\citet{kandasamy16mfbo}.

